\documentclass{compusofia}

\usepackage{biblatex}
\addbibresource{main.bib} % Añada sus referencias en formato BibTeX.
\graphicspath{ {./imagenes/} }

\renewenvironment{abstract}{{\color{colorprincipal}\textbf{RESUMEN }}}{}

%% Aquí reemplace el texto por el nombre de su artículo
\newcommand{\nombrearticulo}{Nombre de su artículo}

%% Aquí reemplace el texto por su nombre
\newcommand{\nombreautor}{Su nombre}

%% Aquí escriba una breve biografía de usted (aparecera al final del artículo).
\newcommand{\descripcionautor}{Su breve biografía.}

\title{\color{colorprincipal} \nombrearticulo}
\author{\nombreautor}
\rhead{\textbf{\nombreautor}}
\lhead{\textbf{\nombrearticulo}}
\lfoot{\color{colorprincipal}\textbf{Compusofía}}
\rfoot{página \thepage}
\providecommand{\keywords}[1]{\textbf{\color{colorprincipal}\textit{Palabras claves:}} #1}

\begin{document}

\maketitle

% Aquí puede reemplazar el texto con el resumen de su artículo.
% No se olvide añadir las palabras claves según las permitidas en el documento de pauta.
\begin{abstract}
Plantilla básica para escribir un artículo en la revista  Compusofía\cite{DBLP:books/sp/Turner18}\footnote{Un ejemplo para añadir una referencia. Y, a su vez, para añadir una nota a pie de página.}.
\end{abstract} \par
\keywords{clave1, clave2}

\section{Introducción}

\section{Código}
A continuación se presenta un ejemplo para añadir código\footnote{El paquete usado es Listing, más información en: \url{https://www.overleaf.com/learn/latex/Code_listing}.}:

%Código
\begin{lstlisting}[language=Python]
def merge_list(left, right):
    if not left or not right:
        return left + right
    else:
        fa = left[0]
        fb = right[0]
    if fa < fb:
        return [fa] + merge_list(left[1:], right)
    else:
        return [fb] + merge_list(left, right[1:])
    
def msort(lst):
    if len(lst) <= 1:
        return lst
    else:
        middle = round(len(lst) / 2)
        left = lst[:middle]
        right = lst[-(len(lst) - middle):]
        return merge_list(msort(left), msort(right))
    
print(msort([9, 2, 1, 8, 3]))
# [1, 2, 3, 8, 9]
\end{lstlisting}

\section{Otros elementos}
En este apartado se presenta un ejemplo de una tabla y de una figura.

%Tabla
\begin{table}[H]
    \begin{center}
        \begin{tabular}{| r | c |}\hline
            Algoritmo & Complejidad (peor caso) \\ \hline
            quicksort &  $O(n^2)$ \\
            merge sort &  $O(n\thinspace log\thinspace n)$  \\
            selection sort & $O(n^2)$ \\ \hline
            \end{tabular}
            \caption{Algoritmos de ordenamiento.}
        \label{tab:algoritmos}
    \end{center}
\end{table}

% Imagen
\begin{figure}[H]
\includegraphics[scale=1]{quicksort}
\centering

\caption{Algoritmo Quicksort.}

\end{figure}

% Aquí van las secciones de su artículo.

\section{Conclusión}
Fin.

\printbibliography

\section*{Sobre el autor}
\begin{description}
    \item[\nombreautor] \descripcionautor  
\end{description}
  
\end{document}
